% Discussion générale en cinq mouvements (trame, étape 10) : rappel, comparaison, explication, contributions, limites et ouverture.
\chapter{Discussion générale}\label{ch:discussion}

\section{Rappel des résultats}
Au premier essai et à conditions égales, les agents Search-R1 battent les harnais fixes (chapitre~\ref{ch:partie1}). Face à leur propre modèle de départ, ils restent meilleurs au premier essai sur les questions de pont, de 4 à 7 points, mais pas sur les questions de comparaison ; et à 32 essais, ils ne résolvent pas plus de questions que lui (chapitre~\ref{ch:partie2}). Selon les règles fixées avant les données, le renforcement a fiabilisé sans élargir ; les données excluent un élargissement de plus de 5 points de pass@32, pas un élargissement de 2 à 4 points. La frontière garde sa taille mais change de composition : l'agent PPO gagne une cinquantaine de questions et en perd une quarantaine, plus que le hasard n'en ferait changer (analyse exploratoire). Trois mesures indépendantes situent le gain : il vient d'abord de la recherche ; l'agent lit aussi mieux, mais ce gain n'apparaît que si on l'oblige à répondre.

\section{Comparaison avec la littérature}\label{sec:discussion-generale}
Le résultat principal transpose aux agents de recherche le constat de \citet{yue2025does} sur les modèles de raisonnement : le renforcement standard rend plus fiable sans rendre possible ce qui ne l'était pas. Il s'accorde aussi avec \citet{yang2026crossovers} : avec une inférence appariée et des bandes simultanées, l'écart à grand nombre d'essais n'est plus distinguable de zéro. Il nuance \citet{zhai2026does}, qui observaient un léger élargissement sur les questions de pont dans un environnement très proche du nôtre. Plusieurs différences peuvent l'expliquer : leur agent était entraîné par eux sur 200 questions seulement, leur modèle faisait 7 milliards de paramètres, leurs tirages étaient au nombre de 64, et leur écart (81 contre 77 questions) n'était pas testé dans le texte principal. Nos données ne contredisent pas un effet de cette taille sur une partie des questions : l'analyse exploratoire par difficulté laisse voir un gain sur les questions que le modèle de départ ne résout presque jamais, que le test sur le sous-ensemble difficile doit confirmer ou infirmer. \aremplir{résultat de la session 4 (sous-ensemble difficile) et de la session 3 ($T = 0$)}

Sur le mécanisme, nos résultats divergent de ceux de \citet{zhai2026does}, qui attribuaient le gain à l'usage de l'information trouvée, et de notre propre chapitre~\ref{ch:partie1}, qui concluait de même. La différence tient au dispositif. La partie~1 comparait des systèmes issus de modèles de départ différents, et sa décomposition reposait sur un indicateur indirect (la réponse apparaît-elle dans un passage lu ?), sans expérience. La partie~2 compare des modèles de même départ, mesure la recherche avec les paragraphes utiles désignés par HotpotQA, et la sépare de la lecture par une manipulation : chaque modèle lit les recherches de l'autre. À modèle de départ identique, la recherche l'emporte. Une partie de « l'avantage de lecture » de la partie~1 venait donc probablement de la différence entre modèle de base entraîné et modèle Instruct, et non du renforcement lui-même.

\section{Explication}
Search-R1 est entraîné avec une récompense de résultat (la réponse finale est-elle exacte ?) et une pénalité qui le retient près de son modèle de référence. Un tel entraînement renforce les trajectoires qui réussissaient déjà parfois : il concentre la probabilité sur elles, sans pousser à explorer celles qui ne réussissaient jamais. C'est le mécanisme de « fiabilisation » que décrit la littérature sans outil, et l'analyse par difficulté en montre la signature : l'agent gagne là où son modèle de départ hésite (1 à 8 réussites sur 16) et ne gagne presque rien là où il échoue toujours. Le tableau n'est pourtant pas celui d'une simple concentration de probabilité, qui laisserait la frontière inchangée question par question : l'agent atteint des questions nouvelles, presque toujours en trouvant les deux paragraphes utiles, et en abandonne d'autres, où sa façon de chercher le mène à une réponse fausse. Le renforcement a déplacé la politique de recherche ; à 3 milliards de paramètres et sur ces questions, ce déplacement s'est soldé par un échange plutôt que par un gain net.

Pourquoi la recherche plutôt que la lecture ? Sur une question de pont, la récompense dépend d'abord du fait de trouver le second paragraphe, dont le titre n'apparaît pas dans la question. Le modèle de départ cherche moins (76~\% de ses trajectoires contiennent une recherche, contre 79 à 94~\% pour les agents) et reformule moins bien ; c'est là que le signal de récompense est le plus fort, et c'est là que l'agent progresse. Sur une question de comparaison, les deux entités sont nommées dans la question : il y a peu à apprendre sur la recherche, et l'agent ne gagne rien.

Enfin, l'agent paie un « impôt » : il ne répond pas dans environ 20~\% de ses trajectoires, contre 2~\% pour son modèle de départ. Pour l'agent PPO, c'est surtout le budget de quatre recherches qui s'épuise ; pour l'agent GRPO, ce sont des sorties qui dérivent et perdent le format. Cet impôt est le prix d'une politique qui a appris à chercher davantage. Il explique aussi le comportement observé en lecture croisée : privé de son raisonnement écrit, l'agent PPO veut chercher encore dans 95~\% des cas, comme si la décision de s'arrêter était portée par ce raisonnement plutôt que par les passages lus.

\section{Contributions}
\textbf{Méthodologique.} Le mémoire applique aux agents de recherche une inférence appariée sur les courbes pass@k (bootstrap à multiplicateurs, bandes simultanées, test d'intersection-union), avec un plan gelé avant les données et deux pièges documentés : le statut d'une trajectoire mal classé par notre propre code, et la régression vers la moyenne quand on classe et mesure sur les mêmes tirages. Il montre aussi qu'un gain de frontière à nombre d'essais fini peut venir de questions rares plutôt que d'un vrai élargissement, d'où la présélection sur des tirages indépendants.

\textbf{Empirique.} Première mesure, à notre connaissance, de la frontière d'agents de recherche publics face à leur propre modèle de départ, et première séparation expérimentale de la recherche et de la lecture pour ces agents.

\textbf{Pratique.} Pour une équipe qui envisage d'entraîner un agent de recherche par renforcement, trois conséquences. Le gain attendu porte sur la fiabilité au premier essai, surtout pour les questions à plusieurs étapes : il faut le mesurer à plusieurs essais, et pas seulement au premier, avant de conclure à un agent « plus capable ». Le taux de sorties sans réponse doit être suivi, car il peut annuler une partie du gain. Enfin, puisque l'essentiel du gain tient à des trajectoires que le modèle de départ produisait déjà parfois, un entraînement supervisé sur les bonnes traces mérite d'être essayé d'abord ; c'est l'hypothèse H4, que ce mémoire n'a pas testée.

\section{Limites et perspectives}
Les résultats portent sur une taille de modèle (3 milliards de paramètres), un jeu de questions (HotpotQA) et un moteur fermé de dix paragraphes par question ; ils ne se transportent pas sans réserve au web ouvert ni aux grands modèles. Les questions de validation de HotpotQA sont de la même distribution que celles de l'entraînement de Search-R1. Le budget de quatre recherches bride l'agent PPO dans 14 à 24~\% de ses trajectoires : avec un budget plus large, sa frontière pourrait bouger, et c'est la première sensibilité à faire. La variante « réponse forcée » de la lecture croisée a été ajoutée après avoir vu les données ; elle est exploratoire, comme les analyses par difficulté et par paragraphes utiles. Trois prolongements suivent directement : tester l'élargissement sur le sous-ensemble difficile et sans recherche (sessions 3 et 4), refaire l'étude avec un budget de recherche plus large, et tester H4, un entraînement supervisé sur les bonnes traces de l'agent, qui répondrait à la question économique posée en introduction.
